\documentclass[11pt,a4paper]{article}
\usepackage[utf8]{inputenc}
\usepackage[french]{babel}
\usepackage[T1]{fontenc}
\usepackage{geometry}
\geometry{margin=2.5cm}
\usepackage{booktabs}
\usepackage{array}
\usepackage{xcolor}
\usepackage{tikz}
\usetikzlibrary{arrows.meta, positioning}
\usepackage{enumitem}
\usepackage{fancyhdr}
\usepackage{titlesec}
\usepackage{parskip}

\definecolor{codeblue}{RGB}{0,90,180}

\pagestyle{fancy}
\fancyhf{}
\rhead{\small Système de Tri Automatique de Bouteilles}
\lhead{\small Rapport Technique}
\cfoot{\thepage}

\titleformat{\section}{\large\bfseries\color{codeblue}}{\thesection}{1em}{}[\titlerule]
\titleformat{\subsection}{\normalsize\bfseries}{\thesubsection}{1em}{}

\begin{document}

% ============================================================
\begin{titlepage}
  \centering
  \vspace*{3cm}
  {\Huge\bfseries Système de Tri Automatique\\[0.5em]
  de Bouteilles PET\par}
  \vspace{1.5cm}
  {\large Rapport Technique\par}
  \vspace{0.5cm}
  \rule{10cm}{0.4pt}\\[0.8em]
  {\normalsize Architecture logicielle et protocoles de communication\par}
  \vspace{1cm}
  {\small Date : \today}
  \vfill
\end{titlepage}

\tableofcontents
\newpage

% ============================================================
\section{Vue d'ensemble du système}

Le système est composé de quatre composants qui communiquent en chaîne :

\vspace{0.8em}
\begin{center}
\begin{tikzpicture}[
  node distance=1.6cm,
  box/.style={rectangle, rounded corners, draw=codeblue, fill=blue!7,
              minimum width=3.2cm, minimum height=1cm, align=center, font=\small},
  arr/.style={->, >=Stealth, thick, draw=gray},
]
  \node[box] (py)    {Script Python\\(détection CV)};
  \node[box, right=2cm of py]   (back)  {Backend\\Node.js};
  \node[box, right=2cm of back] (front) {Frontend\\React};
  \node[box, below=1.5cm of py] (ard)   {Arduino\\Mega 2560};

  \draw[arr] (py)   -- node[above,font=\tiny\itshape]{JSON / CSV} (back);
  \draw[arr] (back) -- node[above,font=\tiny\itshape]{HTTP REST}  (front);
  \draw[arr, bend right=20] (py)  -- node[right,font=\tiny\itshape]{ACC / REJ}   (ard);
  \draw[arr, bend left=20]  (ard) -- node[left,font=\tiny\itshape]{``Finished''} (py);
\end{tikzpicture}
\end{center}

\vspace{0.5em}
\begin{center}
\begin{tabular}{llll}
\toprule
\textbf{Composant} & \textbf{Technologie} & \textbf{Rôle principal} & \textbf{Adresse} \\
\midrule
Script Python  & Python 3 + YOLO + Keras & Détection et décision  & --      \\
Arduino        & C++ (Mega 2560)          & Actionneurs mécaniques & COM6    \\
Backend API    & Node.js + Express        & Pont de données        & :5000   \\
Frontend       & React + TypeScript + Vite& Interface utilisateur  & :8080   \\
\bottomrule
\end{tabular}
\end{center}

% ============================================================
\section{Script Python -- Détection et Décision}

\subsection{Traitement vidéo}

Le script analyse en continu le flux webcam (1280$\times$720, 30 fps) dans un
\textbf{thread GPU dédié} pour ne pas bloquer l'affichage. Chaque frame suit le
pipeline suivant :

\begin{enumerate}[noitemsep]
  \item \textbf{Passe 1 (YOLO)} : détection et tracking des bouteilles sur la frame complète.
  \item \textbf{Passe 2 (YOLO)} : détection du bouchon sur un crop de la bouteille avec un
        padding de 20 px en haut et sur les côtés.
  \item \textbf{CNN ResNet50V2} : classification Remplie / Vide sur le crop 224$\times$224.
  \item \textbf{Lissage} : vote majoritaire sur les 15 dernières frames pour chaque label.
  \item \textbf{Décision} : envoi de la commande à l'Arduino dès que les deux labels sont verrouillés.
\end{enumerate}

\subsection{Modèles d'intelligence artificielle}

\begin{center}
\begin{tabular}{p{3cm} p{5.5cm} p{4.5cm}}
\toprule
\textbf{Modèle} & \textbf{Rôle} & \textbf{Classes détectées} \\
\midrule
YOLOv8         & Détection bouteille et bouchon    & Bouteille, Avec bouchon, Sans bouchon \\
CNN ResNet50V2 & Classification remplissage        & Remplie, Vide \\
\bottomrule
\end{tabular}
\end{center}

\subsection{Lissage et critère de validation}

Pour éviter les fausses détections, chaque label est soumis à un \textbf{vote sur
15 frames} avant d'être verrouillé définitivement :

\begin{itemize}[noitemsep]
  \item \textbf{Remplissage} : bouteille déclarée REMPLIE si $\geq 5$ votes sur 15.
  \item \textbf{Bouchon} : déclaré PRÉSENT si $\geq 8$ votes sur 15.
\end{itemize}

Une bouteille est \textbf{acceptée} uniquement si elle est \textit{vide} \textbf{et}
\textit{sans bouchon}. Dans le cas contraire, elle est rejetée.

% ============================================================
\section{Liaison Arduino}

\subsection{Connexion série}

La connexion est établie automatiquement au démarrage du script. Le programme
parcourt les ports COM disponibles et identifie l'Arduino via des mots-clés
(\textit{arduino}, \textit{ch340}, \textit{ftdi}). Les ports Bluetooth sont
ignorés. La communication utilise un débit de \textbf{9600 bauds}.

\subsection{Protocole de commande}

\begin{center}
\begin{tabular}{p{3.5cm} p{3.5cm} p{6cm}}
\toprule
\textbf{Direction} & \textbf{Message} & \textbf{Signification} \\
\midrule
Python $\to$ Arduino & \texttt{ACC} & Bouteille valide : lancer le cycle de broyage \\
Python $\to$ Arduino & \texttt{REJ} & Bouteille refusée : actionner l'éjection \\
Arduino $\to$ Python & \texttt{Finished} & Cycle de broyage terminé \\
\bottomrule
\end{tabular}
\end{center}

\subsection{Suivi du cycle de broyage}

Lors de l'envoi de \texttt{ACC}, un \textbf{thread d'écoute} démarre et attend
la réponse \texttt{Finished} de l'Arduino. La durée du cycle est calculée et
stockée dans le fichier JSON partagé. Une seule opération peut être en cours à
la fois (flag \texttt{ready}).

\subsection{Traçabilité CSV}

À chaque décision, une ligne est enregistrée dans le fichier
\texttt{detection\_unifie.csv} (Google Drive) :

\begin{center}
\begin{tabular}{lllll}
\toprule
Timestamp & ID Bouteille & État remplissage & Bouchon présent & Statut \\
\midrule
2026-04-27 10:30:00 & 3 & VIDE & NON & ACCEPTEE \\
2026-04-27 10:32:15 & 4 & REMPLIE & OUI & REJETEE \\
\bottomrule
\end{tabular}
\end{center}

% ============================================================
\section{Interface Frontend -- Backend}

\subsection{API REST (Node.js / Express)}

Le backend expose trois endpoints principaux sur le port \textbf{5000} :

\begin{center}
\begin{tabular}{p{4cm} p{2cm} p{7cm}}
\toprule
\textbf{Endpoint} & \textbf{Méthode} & \textbf{Données retournées} \\
\midrule
\texttt{/status}           & GET & Statut temps-réel de la bouteille (existe, remplie, bouchon, valide) \\
\texttt{/csv-stats}        & GET & Statistiques agrégées (total, aujourd'hui, poids, CO\textsubscript{2}) + dernière bouteille \\
\texttt{/broyage-finished} & GET & État du broyage (terminé ou non) + durée en ms \\
\bottomrule
\end{tabular}
\end{center}

Le backend lit en temps réel deux fichiers produits par Python :
\texttt{realtime\_status.json} (statut par frame) et \texttt{detection\_unifie.csv}
(historique des bouteilles acceptées).

\subsection{Interface utilisateur React}

L'interface tourne sur le port \textbf{8080} et présente :

\begin{itemize}[noitemsep]
  \item Les \textbf{statistiques} : nombre total de bouteilles, bouteilles du jour,
        poids recyclé, CO\textsubscript{2} économisé.
  \item Le \textbf{niveau de stockage} PET Transparent (jauge sur 100 bouteilles).
  \item Les \textbf{indicateurs de validation} : Plastique PET / Bouteille vide / Sans bouchon.
  \item Le \textbf{timer de broyage} en temps réel (résolution 100 ms).
  \item Une \textbf{notification de résultat} (verte si acceptée, rouge si refusée)
        avec message personnalisé et bouton OK.
\end{itemize}

\subsection{Machine à états et flux temporel}

\begin{center}
\begin{tikzpicture}[
  state/.style={ellipse, draw=codeblue, fill=blue!7,
                minimum width=2.4cm, minimum height=0.9cm,
                align=center, font=\small},
  arr/.style={->, >=Stealth, draw=gray},
]
  \node[state] (idle) {\texttt{idle}};
  \node[state, right=2.5cm of idle] (det) {\texttt{detecting}};
  \node[state, above right=1.2cm and 2cm of det] (val) {\texttt{valid}};
  \node[state, below right=1.2cm and 2cm of det] (inv) {\texttt{invalid}};

  \draw[arr] (idle) -- node[above,font=\tiny\itshape]{Démarrer} (det);
  \draw[arr] (det)  -- node[above left,font=\tiny\itshape]{OK} (val);
  \draw[arr] (det)  -- node[below left,font=\tiny\itshape]{Refus} (inv);
  \draw[arr, bend left=35]  (val) to node[above,font=\tiny\itshape]{OK} (idle);
  \draw[arr, bend right=35] (inv) to node[below,font=\tiny\itshape]{OK} (idle);
\end{tikzpicture}
\end{center}

\begin{enumerate}[noitemsep]
  \item Clic \textbf{Démarrer} : polling de \texttt{/status} toutes les 350 ms.
  \item Bouteille détectée : attente de \textbf{2 secondes} (verrouillage des labels côté Python).
  \item Lecture de \texttt{/csv-stats} pour récupérer la décision finale.
  \item Si \textbf{valide} : timer broyage + polling \texttt{/broyage-finished} toutes les 800 ms.
  \item Si \textbf{refusée} : notification rouge avec raison précise (bouchon / remplie / les deux).
  \item Clic \textbf{OK} : rechargement de la page, retour à l'état initial.
\end{enumerate}

% ============================================================
\section{Synthèse}

\begin{center}
\begin{tabular}{p{4cm} p{9.5cm}}
\toprule
\textbf{Point clé} & \textbf{Description} \\
\midrule
Robustesse          & Vote sur 15 frames + hystérésis CNN éliminent les faux positifs \\
Découplage          & Python produit des fichiers, Node.js les lit, React les affiche
                      -- aucune dépendance directe \\
Traçabilité         & Chaque bouteille est enregistrée avec horodatage dans le CSV \\
Réactivité          & Polling multi-fréquence selon la phase (350 ms / 800 ms / 100 ms) \\
Sécurité mécanique  & Une seule commande Arduino à la fois grâce au flag \texttt{ready} \\
\bottomrule
\end{tabular}
\end{center}

\end{document}


\definecolor{codeblue}{RGB}{0,100,200}
\definecolor{codegray}{RGB}{100,100,100}
\definecolor{codegreen}{RGB}{0,140,0}
\definecolor{codebg}{RGB}{248,248,248}

\lstset{
  backgroundcolor=\color{codebg},
  basicstyle=\ttfamily\small,
  keywordstyle=\color{codeblue}\bfseries,
  commentstyle=\color{codegreen},
  stringstyle=\color{orange},
  frame=single,
  rulecolor=\color{gray!40},
  breaklines=true,
  numbers=left,
  numberstyle=\tiny\color{codegray},
  xleftmargin=1.5em,
}

\pagestyle{fancy}
\fancyhf{}
\rhead{Système de Tri Automatique de Bouteilles}
\lhead{Rapport Technique}
\cfoot{\thepage}

\titleformat{\section}{\large\bfseries\color{codeblue}}{{\thesection}}{1em}{}[\titlerule]
\titleformat{\subsection}{\normalsize\bfseries}{{\thesubsection}}{1em}{}

\begin{document}

% ============================================================
%  PAGE DE TITRE
% ============================================================
\begin{titlepage}
  \centering
  \vspace*{2cm}
  {\Huge\bfseries Rapport Technique\\[0.4em]
  Système de Tri Automatique de Bouteilles\par}
  \vspace{1.5cm}
  {\large Interface Python -- Liaison Arduino\\[0.3em]
         Interface Frontend -- Backend\par}
  \vspace{2cm}
  \rule{10cm}{0.5pt}\\[0.5em]
  {\normalsize Architecture logicielle, flux de données et protocoles de communication\par}
  \vspace{1cm}
  {\small Date : \today}
  \vfill
\end{titlepage}

% ============================================================
%  TABLE DES MATIÈRES
% ============================================================
\tableofcontents
\newpage

% ============================================================
%  1. VUE D'ENSEMBLE
% ============================================================
\section{Vue d'ensemble du système}

Le système est composé de trois couches logicielles indépendantes qui communiquent
via fichiers partagés et requêtes HTTP :

\vspace{0.5em}
\begin{center}
\begin{tikzpicture}[
  node distance=1.5cm,
  box/.style={rectangle, rounded corners, draw=codeblue, fill=blue!8,
              minimum width=3.8cm, minimum height=1cm, align=center, font=\small},
  arr/.style={->, >=Stealth, thick},
]
  \node[box] (py)  {Script Python\\(\texttt{unified\_bottle\_detection.py})};
  \node[box, right=2.5cm of py] (back) {Backend Node.js\\(\texttt{server.js})};
  \node[box, right=2.5cm of back] (front) {Frontend React\\(\texttt{useRecyclingMachine.ts})};
  \node[box, below=1.8cm of py] (ard)  {Arduino Mega 2560\\(COM6 -- 9600 baud)};

  \draw[arr] (py) -- node[above,font=\tiny]{JSON / CSV} (back);
  \draw[arr] (back) -- node[above,font=\tiny]{HTTP REST} (front);
  \draw[arr, bend right=25] (py) -- node[right,font=\tiny]{ACC / REJ} (ard);
  \draw[arr, bend left=25]  (ard) -- node[left,font=\tiny]{``Finished''} (py);
\end{tikzpicture}
\end{center}

\vspace{0.5em}
\begin{center}
\begin{tabular}{llll}
\toprule
\textbf{Composant} & \textbf{Technologie} & \textbf{Rôle} & \textbf{Port/Chemin} \\
\midrule
Python CV       & Python 3 + YOLO + Keras & Détection + décision & -- \\
Arduino         & C++ (Arduino Mega)       & Actionneurs mécaniques & COM6 \\
Backend API     & Node.js + Express        & Pont données          & :5000 \\
Frontend        & React + TypeScript + Vite& Interface utilisateur  & :8080 \\
\bottomrule
\end{tabular}
\end{center}

% ============================================================
%  2. INTERFACE PYTHON
% ============================================================
\section{Interface Python -- Détection et Décision}

\subsection{Architecture de traitement vidéo}

Le script \texttt{unified\_bottle\_detection.py} implémente une détection en
\textbf{double passe} sur chaque frame webcam, exécutée dans un thread GPU dédié
(\texttt{\_detection\_worker}) afin de ne pas bloquer l'affichage.

\begin{center}
\begin{tikzpicture}[
  node distance=0.9cm,
  flowstep/.style={rectangle, rounded corners, draw=gray, fill=gray!10,
               minimum width=5cm, minimum height=0.7cm, align=center, font=\small},
  arr/.style={->, >=Stealth},
]
  \node[flowstep] (cam)   {Capture webcam (1280$\times$720, 30 fps)};
  \node[flowstep, below of=cam] (p1)   {Passe 1 -- YOLO : tracking bouteilles (classe 0)};
  \node[flowstep, below of=p1]  (p2)   {Passe 2 -- YOLO : détection bouchon (crop + padding 20 px)};
  \node[flowstep, below of=p2]  (cnn)  {CNN ResNet50V2 : classification Remplie / Vide (224$\times$224)};
  \node[flowstep, below of=cnn] (lock) {Lissage 15 frames + verrouillage global};
  \node[flowstep, below of=lock](dec)  {Décision ACC / REJ + écriture CSV + JSON};

  \draw[arr] (cam) -- (p1);
  \draw[arr] (p1)  -- (p2);
  \draw[arr] (p2)  -- (cnn);
  \draw[arr] (cnn) -- (lock);
  \draw[arr] (lock)-- (dec);
\end{tikzpicture}
\end{center}

\subsection{Modèles utilisés}

\begin{tabular}{p{3.5cm} p{4.5cm} p{5cm}}
\toprule
\textbf{Modèle} & \textbf{Fichier} & \textbf{Rôle} \\
\midrule
YOLOv8 (best.pt) & \texttt{Bottle-Bottle-Cap-Detection-System-main/best.pt}
  & Détection bouteille + bouchon. 3 classes : Bottle, Avec Bouchon, Sans Bouchon \\
\addlinespace
CNN ResNet50V2 & \texttt{remplie ou non/artifacts/models/best\_resnet50v2.keras}
  & Classification binaire Remplie / Vide. Score sigmoid : $\approx 0$ = remplie, $\approx 1$ = vide \\
\bottomrule
\end{tabular}

\subsection{Lissage temporel et verrouillage}

Pour éviter les oscillations de détection, les deux labels (remplissage et bouchon)
sont soumis à un \textbf{vote majoritaire sur 15 frames consécutives} :

\begin{lstlisting}[language=Python, caption={Vote majoritaire bouchon}]
self.global_cap_history.append(1 if cap_this_bottle else 0)
if len(self.global_cap_history) >= 15:
    self.global_cap_locked_value = sum(self.global_cap_history) >= 8
    self.global_cap_locked = True   # Verrou definitif
\end{lstlisting}

Seuils appliqués :
\begin{itemize}[noitemsep]
  \item \textbf{Remplissage} : REMPLIE si $\geq 5$ votes sur 15 (seuil bas pour sensibilité).
  \item \textbf{Bouchon} : AVEC si $\geq 8$ votes sur 15 (seuil médian).
  \item \textbf{Hystérésis CNN} : seuil 0.50 (VIDE$\to$REMPLIE) / 0.72 (REMPLIE$\to$VIDE).
\end{itemize}

Une fois verrouillés, les labels ne changent plus jusqu'à disparition de la bouteille
(grace period de 8 frames sans détection avant reset).

\subsection{Critère de validation}

\[
  \text{Bouteille acceptée} \iff \texttt{global\_locked\_label} = \text{VIDE}
  \;\wedge\; \texttt{global\_cap\_locked\_value} = \texttt{False}
\]

Si le critère est satisfait : envoi \texttt{ACC\textbackslash n} à l'Arduino.
Sinon : envoi \texttt{REJ\textbackslash n}.

% ============================================================
%  3. LIAISON ARDUINO
% ============================================================
\section{Liaison Arduino}

\subsection{Initialisation de la connexion série}

La connexion est établie dans \texttt{init\_arduino()} :

\begin{lstlisting}[language=Python, caption={Connexion série Arduino}]
self.arduino_serial = serial.Serial(
    port=arduino_port,   # auto-detection parmi les ports COM
    baudrate=9600,
    timeout=1
)
time.sleep(2)  # Attente reset Arduino apres ouverture port
\end{lstlisting}

La détection automatique du port cherche les mots-clés :
\texttt{arduino}, \texttt{ch340}, \texttt{usb serial}, \texttt{ftdi}.
Les ports Bluetooth sont ignorés.

\subsection{Protocole de communication}

\begin{center}
\begin{tabular}{p{3cm} p{4cm} p{6cm}}
\toprule
\textbf{Direction} & \textbf{Message} & \textbf{Signification} \\
\midrule
Python $\to$ Arduino & \texttt{ACC\textbackslash n} & Bouteille valide : démarrer le cycle de broyage \\
Python $\to$ Arduino & \texttt{REJ\textbackslash n} & Bouteille refusée : éjection \\
Arduino $\to$ Python & \texttt{Finished}            & Cycle de broyage terminé \\
\bottomrule
\end{tabular}
\end{center}

\subsection{Suivi de la durée de broyage}

\begin{lstlisting}[language=Python, caption={Tracking du cycle de broyage}]
# Envoi ACC + demarrage chronometre
self.arduino_serial.write(b'ACC\n')
self.broyage_start_time = time.time()
self.broyage_finished = False

# Thread en attente de "Finished"
threading.Thread(target=self.wait_for_finished, daemon=True).start()

# Dans wait_for_finished() :
if feedback == "Finished":
    self.broyage_duration_ms = round(
        (time.time() - self.broyage_start_time) * 1000
    )
    self.broyage_finished = True
    self.ready = True
\end{lstlisting}

Les champs \texttt{broyage\_finished} et \texttt{broyage\_duration\_ms} sont
écrits dans \texttt{data\_logs/realtime\_status.json} à chaque frame.

\subsection{Écriture CSV}

À chaque bouteille avec labels verrouillés, une ligne est ajoutée au fichier
\texttt{G:\textbackslash Mon Drive\textbackslash csv file robotique\textbackslash detection\_unifie.csv} :

\begin{lstlisting}[caption={En-tête CSV}]
Timestamp, ID_Bouteille, Etat_Remplissage, Bouteille_Avec_Bouchon, Statut
\end{lstlisting}

\texttt{Statut} vaut \texttt{ACCEPTEE} ou \texttt{REJETEE}.

% ============================================================
%  4. INTERFACE FRONTEND -- BACKEND
% ============================================================
\section{Interface Frontend -- Backend}

\subsection{API REST Backend (Node.js / Express)}

Le backend tourne sur \texttt{http://localhost:5000} et expose 4 endpoints :

\begin{center}
\begin{tabular}{p{3.5cm} p{2cm} p{7.5cm}}
\toprule
\textbf{Endpoint} & \textbf{Méthode} & \textbf{Description} \\
\midrule
\texttt{/status}           & GET & Lit \texttt{realtime\_status.json} et retourne le statut
                                    temps-réel de la bouteille. \\
\addlinespace
\texttt{/csv-stats}        & GET & Lit \texttt{detection\_unifie.csv} et retourne les
                                    statistiques agrégées + \texttt{lastBottle}. \\
\addlinespace
\texttt{/broyage-finished} & GET & Retourne \texttt{\{finished, durationMs\}} depuis le JSON. \\
\addlinespace
\texttt{/detect}           & POST & (non utilisé en production) Spawn du script Python. \\
\bottomrule
\end{tabular}
\end{center}

\subsubsection*{Réponse type de \texttt{/status}}
\begin{lstlisting}[language=JavaScript]
{
  "bottleExists": true,
  "isFilled": false,
  "hasCap": false,
  "valid": true,
  "rejectionReason": null,
  "numBottles": 1
}
\end{lstlisting}

\subsubsection*{Réponse type de \texttt{/csv-stats}}
\begin{lstlisting}[language=JavaScript]
{
  "totalBottles": 42,  "todayBottles": 5,
  "totalWeight": 1.26, "co2Saved": 3.15,
  "transparentCount": 42, "coloredCount": 0,
  "lastBottle": { "isFilled": false, "hasCap": false,
                  "valid": true, "timestamp": "2026-04-27 10:30:00" }
}
\end{lstlisting}

\subsection{Hook React \texttt{useRecyclingMachine}}

Le hook centralise toute la logique de l'interface utilisateur. Il gère une
machine à états avec les états suivants :

\begin{center}
\begin{tikzpicture}[
  state/.style={ellipse, draw=codeblue, fill=blue!8,
                minimum width=2.2cm, minimum height=0.8cm,
                align=center, font=\small},
  arr/.style={->, >=Stealth, font=\tiny},
]
  \node[state] (idle)      {\texttt{idle}};
  \node[state, right=2cm of idle] (det) {\texttt{detecting}};
  \node[state, above right=1cm and 2cm of det] (val) {\texttt{valid}};
  \node[state, below right=1cm and 2cm of det] (inv) {\texttt{invalid}};

  \draw[arr] (idle) -- node[above]{clic ``Démarrer''} (det);
  \draw[arr] (det)  -- node[above left]{bouteille OK} (val);
  \draw[arr] (det)  -- node[below left]{refusée/absente} (inv);
  \draw[arr, bend left=40] (val) to node[above]{OK (reload)} (idle);
  \draw[arr, bend right=40] (inv) to node[below]{OK (reload)} (idle);
\end{tikzpicture}
\end{center}

\subsection{Flux temporel complet}

\begin{enumerate}[noitemsep]
  \item Clic \textbf{Démarrer} $\to$ \texttt{status = 'detecting'} + démarrage polling
        \texttt{/status} toutes les \textbf{350 ms}.
  \item Dès que \texttt{bottleExists = true} $\to$ \textbf{attente 2 secondes}
        (Python verrouille les labels et écrit le CSV).
  \item Appel \texttt{/csv-stats} : récupération \texttt{lastBottle} + mise à jour stats.
  \item \textbf{Si valide} : démarrage timer 100 ms (affichage chrono)
        + polling \texttt{/broyage-finished} toutes les \textbf{800 ms}.
  \item Dès que \texttt{finished = true} : arrêt timers, affichage notification verte
        avec durée + bouton OK.
  \item \textbf{Si refusée} : \texttt{status = 'invalid'}, affichage notification rouge
        personnalisée (bouchon / remplie / les deux) + bouton OK.
  \item Clic OK $\to$ \texttt{window.location.reload()}.
\end{enumerate}

\subsection{Paramètres de timing}

\begin{center}
\begin{tabular}{ll}
\toprule
\textbf{Paramètre} & \textbf{Valeur} \\
\midrule
Délai avant lecture CSV      & 2\,000 ms \\
Polling \texttt{/status}     & 350 ms \\
Polling \texttt{/broyage-finished} & 800 ms \\
Timer chrono broyage         & 100 ms \\
Timeout global détection     & 15\,000 ms \\
\bottomrule
\end{tabular}
\end{center}

% ============================================================
%  5. FICHIERS D'ÉCHANGE
% ============================================================
\section{Fichiers d'échange entre les couches}

\begin{center}
\begin{tabular}{p{5.5cm} p{2.5cm} p{6cm}}
\toprule
\textbf{Fichier} & \textbf{Format} & \textbf{Contenu} \\
\midrule
\texttt{data\_logs/realtime\_status.json}
  & JSON & Statut temps-réel : \texttt{bouteille\_existe},
    \texttt{bouteille\_remplie}, \texttt{bouteille\_avec\_bouchon},
    \texttt{broyage\_finished}, \texttt{broyage\_duration\_ms} \\
\addlinespace
\texttt{G:\textbackslash{}Mon Drive\textbackslash{}...\quad detection\_unifie.csv}
  & CSV & Historique des bouteilles : horodatage, ID, remplissage,
    bouchon, statut (ACCEPTEE/REJETEE) \\
\bottomrule
\end{tabular}
\end{center}

Le JSON est écrit à \textbf{chaque frame} (environ 30 fps), le CSV uniquement
lors du verrouillage des labels (une ligne par bouteille).

% ============================================================
%  6. RÉSUMÉ
% ============================================================
\section{Résumé}

\begin{center}
\begin{tabular}{p{4cm} p{9.5cm}}
\toprule
\textbf{Point fort} & \textbf{Description} \\
\midrule
Robustesse détection     & Vote 15 frames + hystérésis CNN évitent les faux positifs \\
Découplage des couches   & Python écrit fichiers $\to$ Node.js lit $\to$ React affiche,
                           sans dépendance directe \\
Réactivité UI            & Polling asynchrone multi-fréquence selon la phase du cycle \\
Traçabilité              & CSV persistant sur Google Drive avec horodatage \\
Sécurité mécanique       & Arduino ne reçoit une commande que si \texttt{ready = True}
                           (une seule opération à la fois) \\
\bottomrule
\end{tabular}
\end{center}

\end{document}
